% Junior Level 3, Week 1 -- answer key
\answersweek{1}{Patterns in Patterns}{Answers for workbook chapter 1 and play-at-home card 1.}

\section{Lesson questions}
\begin{answer}[Exit question]
4 squares. The inside repeat draws one square; the outside repeat runs it 4 times, with one quarter turn after each square, so the four squares meet at the middle like a window. A child who says \enquote{16 lines} has counted every move and is also right about the lines.
\end{answer}

\begin{answer}[Reflection]
\enquote{A repeat inside a repeat} or \enquote{a pattern inside a pattern} both show understanding; pointing at the two orange repeats in their script also counts. Any drawing of a pattern is right: ask which part repeats.
\end{answer}

\section{Workbook exercises}
\begin{answer}[Exercise 1.1: The Window]
\textbf{What you should see:} after the green flag, a small cat facing right draws four squares one after the other: first the one to the lower right, then lower left, upper left and upper right. Together they look like one big square window with four panes. The cat ends in the middle, facing right.

\textbf{What to look for:} \tickbox{} one repeat sits inside the other \quad \tickbox{} the last turn is under the inside repeat but inside the outside repeat \quad \tickbox{} \enquote{pen down} is above the repeats.
\end{answer}

\begin{answer}[Exercise 1.2: Square Flower]
\textbf{What you should see:} six squares that overlap and share one corner in the middle, each turned a little more than the one before, like the petals of a flower. The cat ends where it started.

\textbf{What to look for:} \tickbox{} a pink-red \enquote{define square} block with the repeat 4 under it \quad \tickbox{} the flag script has repeat 6 with \enquote{square} and \enquote{turn 60} inside \quad \tickbox{} the 60 came from the turn card. A plain repeat 4 inside the repeat 6 (no My Block) draws the same flower and is fine on the support path.
\end{answer}

\begin{answer}[Exercise 1.3: The One-Square Flower]
\textbf{The bug:} circle the \enquote{repeat 4} at the bottom. It sits \textbf{under} the \enquote{repeat 6}, so it runs only once, after the cat has finished spinning. With the bug the cat turns round on the spot six times (no line, because it does not move), then draws one square.

\textbf{Fix:} drag the \enquote{repeat 4} into the \enquote{repeat 6}, above the turn. Now each time round, the cat draws a square and then turns 60: six squares, the same flower as 1.2.
\end{answer}

\section{Play-at-home card}
\begin{answer}[Pattern maker]
For \enquote{(clap, clap, pat) 3 times, then jump; all of that 2 times} there are 12 claps, 6 pats and 2 jumps. Any pattern with a small part inside a bigger part is right. Look for the child checking whether the grown-up did it right.
\end{answer}

\begin{answer}[Draw a nested pattern]
Two groups, each with three rows of circle, circle, triangle and then one star: 12 circles, 6 triangles and 2 stars in all. Any colors are right.
\end{answer}
